\scanpage{166}
\mspage{40}

En effet, on a

\uline{Proposition} 4.4.2. Soit $\unsure{\mathcal M|M}$ stable par \inserted{petites} $\unsure{\underrightarrow{\lim}|\underline{\lim}}$ inductives
filtrantes, et $W$ \inserted{$\subset\mathrm{Fl}\mathcal M$} un localiseur \inserted{(mult. saturé)}
\uline{stable par les dites limites}. Alors $\mathrm{Cof}_W$ satisfait les
conditions de stabilité (4.4.3) (``stabilité
de Quillen gauche''), et est \unsure{stable} aussi
par $\underrightarrow{\lim}$ filtrantes quelconques.

En effet, la condition a) est satisfaite
sans condition \unsure{suppl.} sur $W$. D'autre part,
il est immédiat que cette condition
\unsure{p.r|pour|par} aux $\underrightarrow{\lim}$ filtrantes implique que
$\mathrm{Cof}_W$ \uline{est stable par limites inductives
filtrantes} dans $\uline{\mathrm{Fl}}\,\mathcal M$.

\uline{Cor. 4.4.3} Sous les conditions de la prop.
4.4.\unsure{2|7}., $\mathrm{Cof}_W\cap W=W^{\mathrm{univ}}$ (\unsure{Ens.|ens.} des
$W$-équivalences co-universelles dans $\mathcal M$) est
stable par limites inductives filtrantes,
et \deleted{est donc est} donc ``Qstable à gauche''
au sens de (4.4.3.).

\uline{Corollaire}. 4.4.4. Sous les conditions de
% slanted
la prop. 4.4.2\deleted{\unsure{,}}*, \inserted{supposons de plus que les
sommes amalgamées finies existent dans
$\mathcal M$, et} \deleted{supposons de plus que} soit $C_0$ une \inserted{petite} partie de $\mathrm{Cof}_W$,
\inserted{dont les éléments soient accessibles, \deleted{et \unsure{coaccessibles}}}
\deleted{$\mathrm{Fl}\mathcal M$ \unsure{close|clos} à gauche. Alors $\unsure{C\cap W}$ satisfait}
\inserted{\unsure{et} considérons}
\[
 C=\widetilde C_0,\quad TC\overset{\mathrm{déf}}{=}C\cap W.
\]
Alors on a
\[
 \deleted{\unsure{F}}C\subset\mathrm{Cof}_W,\quad TC=C\cap W^{\mathrm{univ}}\subset W^{\mathrm{univ}},
\]

% archive sheet number: 165
